\subsection{PROPERTIES OF E ⊗\_A F RELATIVE TO EXACT SEQUENCES}

PROPOSITION 5. Let E, E', E'' be right A-modules, F a left A-module and

\[
\mathrm{E} ^ {\prime} \stackrel {u} {\longrightarrow} \mathrm{E} \stackrel {v} {\longrightarrow} \mathrm{E} ^ {\prime \prime} \longrightarrow 0\tag{14}
\]

an exact sequence of linear mappings. Writing \(\bar{u} = u \otimes 1_{\mathbb{F}}, \bar{v} = v \otimes 1_{\mathbb{F}}\) , the sequence

\[
\mathrm{E} ^ {\prime} \otimes_ {\mathrm{A}} \mathrm{F} \xrightarrow {\bar {u}} \mathrm{E} \otimes_ {\mathrm{A}} \mathrm{F} \xrightarrow {\bar {v}} \mathrm{E} ^ {\prime \prime} \otimes_ {\mathrm{A}} \mathrm{F} \longrightarrow 0\tag{15}
\]

of Z-homomorphisms is exact.

By virtue of no. 2, formula (5), \(\bar{v} \circ \bar{u} = (v \circ u) \otimes 1_{\mathbf{F}} = 0\) ; the image \(\mathrm{H} = \bar{u} (\mathbf{E}' \otimes \mathbf{F})\) is contained in the kernel \(\mathbf{L} = \operatorname{Ker}(\bar{v})\) ; by passing to the quotient, we therefore derive from \(\bar{v}\) a \(\mathbf{Z}\) -linear mapping \(f\) of the cokernel \(\mathbf{M} = (\mathbf{E} \otimes \mathbf{F}) / \mathrm{H}\) of \(\bar{u}\) into \(\mathbf{E}'' \otimes \mathbf{F}\) ; it must be proved that \(f\) is bijective and it will hence suffice to define a \(\mathbf{Z}\) -linear mapping \(g: \mathbf{E}'' \otimes \mathbf{F} \to \mathbf{M}\) such that \(g \circ f\) and \(f \circ g\) are the identity mappings.

Let \(x'' \in E'', y \in F\) ; by hypothesis there exists \(x \in E\) such that \(v(x) = x''\) . We show that, if \(x_1, x_2\) are two elements of \(E\) such that \(v(x_1) = v(x_2) = x''\) and \(\phi: E \otimes F \to M\) is the canonical mapping, then \(\phi(x_1 \otimes y) = \phi(x_2 \otimes y)\) . It suffices to prove that if \(v(x) = 0\) then \(\phi(x \otimes y) = 0\) , which follows from the fact that \(x = u(x')\) with \(x' \in E'\) , whence \(x \otimes y = u(x') \otimes y = \bar{u}(x' \otimes y) \in H\) . If \((x'', y)\) is mapped to the unique value of \(\phi(x \otimes y)\) for all \(x \in E\) such that \(v(x) = x''\) , a mapping is defined of \(E'' \times F\) into \(M\) ; this mapping is Z-bilinear and satisfies conditions (1) (no. 1), since \(v(x\lambda) = x''\lambda\) and \((x\lambda) \otimes y = x \otimes (\lambda y)\) for \(x \in E\) ; hence there is a Z-linear mapping \(g\) of \(E'' \otimes F\) into \(M\) such that \(g(x'' \otimes y) = \phi(x \otimes y)\) for \(y \in F, x \in E\) and \(x'' = v(x)\) . This definition further proves that \(f \circ g\) coincides with the identity mapping for the elements of

\(E'' \otimes F\) of the form \(x'' \otimes y\) and hence \(f \circ g\) is the identity mapping of \(E'' \otimes F\) ; on the other hand, for \(x \in E\) and \(y \in F\) , \(f(\phi(x \otimes y)) = v(x) \otimes y\) by definition, hence \(g(f(\phi(x \otimes y))) = \phi(x \otimes y)\) and, as the elements of the form \(\phi(x \otimes y)\) generate M, \(g \circ f\) is the identity mapping of M.

COROLLARY. Let F, F', F'' be left A-modules, E a right A-module and

\[
\mathrm{F} ^ {\prime} \stackrel {s} {\longrightarrow} \mathrm{F} \stackrel {t} {\longrightarrow} \mathrm{F} ^ {\prime \prime} \longrightarrow 0\tag{16}
\]

an exact sequence of linear mappings. Writing \(\bar{s} = 1_{\mathbb{E}} \otimes s\) , \(\bar{t} = 1_{\mathbb{E}} \otimes t\) , the sequence of \(\mathbf{Z}\) -homomorphisms

\[
\mathrm{E} \otimes_ {\mathrm{A}} \mathrm{F} ^ {\prime} \stackrel {\hat {s}} {\longrightarrow} \mathrm{E} \otimes_ {\mathrm{A}} \mathrm{F} \stackrel {\hat {t}} {\longrightarrow} \mathrm{E} \otimes_ {\mathrm{A}} \mathrm{F} ^ {\prime \prime} \longrightarrow 0\tag{17}
\]

is exact.

When E (resp. F) is considered as a left (resp. right) A \(^{0}\) -module, F ⊗A \(^{0}\) E is identified with E ⊗A F and there are analogous identifications for F' ⊗A \(^{0}\) E and F'\,' ⊗A \(^{0}\) E (no. 1, Corollary 2 to Proposition 1); the corollary then follows immediately from Proposition 5.

Remark. Note that in general, if \(E'\) is a submodule of a right A-module E and \(j:E'\to E\) the canonical injection, the canonical mapping

\[
j \otimes 1 _ {\mathbf {F}}: \mathbf {E} ^ {\prime} \otimes \mathbf {F} \rightarrow \mathbf {E} \otimes \mathbf {F}
\]

is not necessarily injective. In other words, for an exact sequence

\[
0 \longrightarrow \mathrm{E} ^ {\prime} \stackrel {u} {\longrightarrow} \mathrm{E} \stackrel {v} {\longrightarrow} \mathrm{E} ^ {\prime \prime} \longrightarrow 0\tag{18}
\]

it cannot in general be concluded that the sequence

\[
0 \longrightarrow \mathrm{E} ^ {\prime} \otimes \mathrm{F} \stackrel {u} {\longrightarrow} \mathrm{E} \otimes \mathrm{F} \stackrel {v} {\longrightarrow} \mathrm{E} ^ {\prime \prime} \otimes \mathrm{F} \longrightarrow 0\tag{19}
\]

is exact.

Take for example A = Z, E = Z, \(E' = 2Z\) , F = Z/2Z. As \(E'\) is isomorphic to E, \(E' \otimes F\) is isomorphic to E \(\otimes\) F which is itself isomorphic to F (no. 4, Proposition 4). But for all \(x' = 2x \in E'\) (where \(x \in E\) ) and all \(y \in F\) , \(j(x') \otimes y = (2x) \otimes y = x \otimes (2y) = 0\) , since 2y = 0, and the canonical image of \(E' \otimes F\) in E \(\otimes\) F reduces to 0.

In other words, care must be taken to distinguish, for a submodule \(E'\) of E and an element \(x \in E'\) , between the element \(x \otimes y\) ``calculated in \(E' \otimes F\)'' and the element \(x \otimes y\) ``calculated in \(E \otimes F\)'' (in other words, the element \(j(x) \otimes y\) ).

We shall study later, under the name of flat modules, the modules F such that the sequence (19) is exact for every exact sequence (18) (Commutative Algebra, I, § 2).

PROPOSITION 6. Given two exact sequences (14) and (16), the homomorphism \(v \otimes t: \mathbf{E} \otimes_{\mathbf{A}} \mathbf{F} \to \mathbf{E}'' \otimes_{\mathbf{A}} \mathbf{F}''\) is surjective and its kernel is equal to

\[
\operatorname{Im} (u \otimes \mathrm{l} _ {\mathrm{F}}) + \operatorname{Im} (\mathrm{l} _ {\mathrm{E}} \otimes s)
\]

Now \(v \otimes t = (v \otimes 1_{\mathbb{F}''}) \circ (1_{\mathbb{E}} \otimes t)\) (no. 2, formula (5)) and \(v \otimes t\) is therefore surjective, being the composition of two surjective homomorphisms by virtue of Proposition 5 and its Corollary. On the other hand, for \(z \in \mathbb{E} \otimes \mathbb{F}\) to be in the kernel of \(v \otimes t\) , it is necessary and sufficient that \((1_{\mathbb{E}} \otimes t) (z)\) belong to the kernel of \(v \otimes 1_{\mathbb{F}''}\) , that is, by virtue of (15) to the image of

\[
u \otimes 1 _ {F ^ {\prime \prime}}: \mathrm{E} ^ {\prime} \otimes \mathrm{F} ^ {\prime \prime} \rightarrow \mathrm{E} \otimes \mathrm{F} ^ {\prime \prime}.
\]

But as the homomorphism \(t: \mathbf{F} \to \mathbf{F}''\) is surjective, so is

\[
\mathbf {l} _ {\mathbf {E} ^ {\prime}} \otimes t: \mathbf {E} ^ {\prime} \otimes \mathbf {F} \rightarrow \mathbf {E} ^ {\prime} \otimes \mathbf {F} ^ {\prime \prime}
\]

by the Corollary to Proposition 5, hence the condition on \(z\) reduces to the existence of an \(a \in \mathbf{E}' \otimes \mathbf{F}\) such that

\[
(1 _ {\mathrm{E}} \otimes t) (z) = (u \otimes t) (a).
\]

Let \(b = z - (u \otimes 1_{\mathbb{F}})(a)\) ; then \((1_{\mathbb{E}} \otimes t)(b) = 0\) , and by virtue of (17), \(b\) belongs to the image of \(1_{\mathbb{E}} \otimes s\) , which proves the proposition.

In other words:

COROLLARY 1. Let \(\mathbf{E}'\) be a submodule of a right A-module \(\mathbf{E}\) , \(\mathbf{F}'\) a submodule of a left A-module \(\mathbf{F}\) and \(\operatorname{Im}(\mathbf{E}' \otimes_{\mathbf{A}} \mathbf{F})\) and \(\operatorname{Im}(\mathbf{E} \otimes_{\mathbf{A}} \mathbf{F}')\) the sub- \(\mathbf{Z}\) -modules of \(\mathbf{E} \otimes_{\mathbf{A}} \mathbf{F}\) , the respective images of the canonical mappings \(\mathbf{E}' \otimes_{\mathbf{A}} \mathbf{F} \to \mathbf{E} \otimes_{\mathbf{A}} \mathbf{F}\) ,

\[
\mathbf {E} \otimes_ {\mathbf {A}} \mathbf {F} ^ {\prime} \rightarrow \mathbf {E} \otimes_ {\mathbf {A}} \mathbf {F}.
\]

Then there is a canonical Z-module isomorphism

\[
\pi : (\mathrm{E} / \mathrm{E} ^ {\prime}) \otimes_ {\mathrm{A}} (\mathrm{F} / \mathrm{F} ^ {\prime}) \rightarrow (\mathrm{E} \otimes_ {\mathrm{A}} \mathrm{F}) / (\operatorname{Im} (\mathrm{E} ^ {\prime} \otimes_ {\mathrm{A}} \mathrm{F}) + \operatorname{Im} (\mathrm{E} \otimes_ {\mathrm{A}} \mathrm{F} ^ {\prime}))\tag{20}
\]

such that, for \(\xi \in \mathbf{E} / \mathbf{E}'\) , \(\eta \in \mathbf{F}\) , \(\pi (\xi \otimes \eta)\) is the class of all elements \(x \otimes y \in \mathbf{E} \otimes_{\mathbf{A}} \mathbf{F}\) such that \(x \in \xi\) and \(y \in \eta\) .

Note that when E is a \(((B_{i}^{\prime}); A, (C_{j}^{\prime}))\) -multimodule, F a \((A, (B_{h}^{\prime\prime}); (C_{k}^{\prime\prime}))\) -multimodule and \(E^{\prime}\) and \(F^{\prime}\) submultimodules of E and F respectively, the isomorphism (20) is an isomorphism for the \(((B_{i}^{\prime}), (B_{h}^{\prime\prime}); (C_{j}^{\prime}), (C_{k}^{\prime\prime}))\) -multimodule structures of the two sides (no. 3).

COROLLARY 2. Let \(\mathfrak{a}\) be a right ideal of \(\mathbf{A}\) , \(\mathbf{F}\) a left A-module and \(\mathfrak{a}\mathbf{F}\) the sub- \(\mathbf{Z}\) -module of \(\mathbf{F}\) generated by the elements of the form \(\lambda x\) , where \(\lambda \in \mathfrak{a}\) and \(x \in \mathbf{F}\) . Then there is a canonical \(\mathbf{Z}\) -module isomorphism

\[
\pi : (\mathrm{A} / \mathfrak {a}) \otimes_ {\mathrm{A}} \mathrm{F} \rightarrow \mathrm{F} / \mathfrak {a} \mathrm{F}\tag{21}
\]

such that, for all \(\bar{\lambda} \in \mathbf{A} / \mathfrak{a}\) and all \(x \in \mathbf{F}\) , \(\pi(\bar{\lambda} \otimes x)\) is the class mod. \(\mathfrak{aF}\) of \(\lambda x\) , where \(\lambda \in \bar{\lambda}\) .

In particular, for A = Z, it is seen that for every integer n and every Z-module F, \((\mathbf{Z}/n\mathbf{Z}) \otimes_{\mathbf{Z}} F\) is canonically identified with the quotient Z-module F/nF.

COROLLARY 3. Let A be a commutative ring, a an ideal of A and E and F two A-modules such that a is contained in the annihilator of F. Then the (A/a)-modules E ⊗\_A F and (E/aE) ⊗\_\{A/a\} F are canonically isomorphic.

F and \(E \otimes_{A} F\) are annihilated by a and hence have canonical (A/a)-module structures (§1, no. 12) and if we write \(E' = aE\) , then \(\operatorname{Im}(E' \otimes_{A} F) = 0\) ; then there is a canonical isomorphism (20) of \(E \otimes_{A} F\) onto \((\mathrm{E}/a\mathrm{E}) \otimes_{A} F\) and the latter is itself identical with \((\mathrm{E}/a\mathrm{E}) \otimes_{A/a} F\) (no. 3, Corollary to Proposition 2).

COROLLARY 4. Let \(\mathfrak{a}, \mathfrak{b}\) be two ideals in a commutative ring \(\mathbf{C}\) ; the \(\mathbf{C}\) -module \((\mathbf{C} / \mathfrak{a}) \otimes_{\mathbb{C}} (\mathbf{C} / \mathfrak{b})\) is then canonically isomorphic to \(\mathbf{C}/(\mathfrak{a} + \mathfrak{b})\) .

